% §1035's `wrapup` puts an empty discretionary after an explicit hyphen:
%
%     @d wrapup(#)==if cur_l<non_char then
%       begin if link(cur_q)>null then
%         if character(tail)=qi(hyphen_char[main_f]) then ins_disc:=true;
%       if ligature_present then pack_lig(#);
%       if ins_disc then
%         begin ins_disc:=false;
%         if mode>0 then tail_append(new_disc);
%         end;
%       end
%
% That empty discretionary is the breakpoint after `well-` in `well-known`.
% §869 charges `\exhyphenpenalty` for it rather than `\hyphenpenalty`, because
% its pre-break list is empty — there is no second hyphen to set.
%
% `wrapup` runs when the *character run* ends, so the discretionary lands
% immediately after the hyphen and before whatever ended the run. In
% `length-$a$` the run ends at the `$`, and the node order is
%
%     l e n g t h - \discretionary  <math> a <math>
%
% not `… - <math> a <math> \discretionary`. Putting it after the formula moves
% the breakpoint one atom to the right, which makes the line that ends there
% longer and its badness lower — so a first pass that TeX cannot complete
% within `\pretolerance` succeeded here instead, the paragraph never reached
% the second pass, and a word TeX hyphenates came out whole. 2608.18024v1
% matched pdfTeX's 60 pages and was 14 characters short, all of them hyphens.
%
% `mode>0` in the last line is why this is a *paragraph* and not an `\hbox`:
% restricted horizontal mode gets no such discretionary at all.
\catcode`\{=1 \catcode`\}=2 \catcode`\$=3
% Without this the font's `\hyphenchar` is `\defaulthyphenchar`, which INITEX
% leaves at 0 — §1035 compares the character against it and never fires, so
% every case below shows no discretionary in *either* engine and the file
% asserts nothing. plain.tex sets `\defaulthyphenchar` to `\-`; this file does
% not load plain.
\font\r=cmr10 \r \hyphenchar\r=`\- \hsize=200pt \parindent=0pt
\font\mi=cmmi10 \font\sy=cmsy10 \font\ex=cmex10
\textfont0=\r \scriptfont0=\r \scriptscriptfont0=\r
\textfont1=\mi \scriptfont1=\mi \scriptscriptfont1=\mi
\textfont2=\sy \scriptfont2=\sy \scriptscriptfont2=\sy
\textfont3=\ex \scriptfont3=\ex \scriptscriptfont3=\ex
\showboxdepth=2 \showboxbreadth=2000 \tracingonline=0

% The discretionary has to be looked for in the *unbroken* horizontal list:
% `\showbox` of the finished `\vbox` shows the lines after `post_line_break`,
% and there the placement difference does not survive. `\showlists` in the
% middle of the paragraph is what sees it. (A `\showbox` version of these cases
% passed happily against a build that placed the node two positions late.)
%
% A hyphen with a formula straight after it: the run ends at the `$`, so the
% discretionary goes between the hyphen and `\mathon` —
%
%     h - \discretionary \mathon a \mathoff
%
% not after the formula, and not after the space that follows it.
\noindent length-$a$ submodules\showlists
\par

% A hyphen at the end of an ordinary word, the everyday case: the run ends at
% the space.
\noindent well-known result\showlists
\par

% Two hyphens running, then a formula.
\noindent a--$b$\showlists
\par

% A hyphen with *nothing* after it: no further node is appended, so the only
% thing that can put the discretionary down is the run ending at `\showlists`.
\noindent trailing-\showlists
\par

% In restricted horizontal mode there is no discretionary at all: `mode>0`.
\setbox0=\hbox{length-$a$}
\showbox0
\end
